\subsection{The residual full-group problem and the local kernels}

\begin{statement}{A.2}{Open}{The full profinite procyclic normalizer criterion}
We do not prove or refute that every nontrivial closed procyclic
$A\subset\pi$ with open arithmetic normalizer is cuspidal.
The full fixed-group theorem A.1 does not supply a section of that
normalizer. The residual possibility is contained in
$R_{\rm nil}(\pi')$, and the prime-complement configuration of
LM6 A.62 remains an embedding problem inside the actual arithmetic
curve group. The first-cover results A.11--A.14 rule out the
specified additional pro-prime detections but do not form a cofinal
family of all finite covers.
\end{statement}
\begin{proof}[Attempt and remaining obstruction]
The strengthened finite-cover argument is A.15 below. The structural
reduction has $N_0=P\rtimes B$, geometric kernel $P\times C$,
$P\simeq\Z_\ell$, and $C$ pro-prime-to-$\ell$. It supplies a
section of $N_0/C$, not one of $N_0$. The explicit nonsplit abstract
extensions in LM6 A.63 show why prime support alone does not close
this step. The action on $P$ is constrained by the CM-character
argument of LM6 A.53, but that is a statement about the quotient
action and does not annul the extension class of $B$.

We also tried a direct local valuation argument. At a good place
of residue characteristic different from $\ell$, the preimage in
$B$ of a pro-$\ell$ Sylow subgroup of the local Galois group has a
pro-$\ell$ Sylow complement mapping isomorphically to that subgroup:
its kernel is pro-prime-to-$\ell$. This gives an actual local
section on that Sylow subgroup. Tame inertia acts trivially on the
unramified CM character of $P$. The uniqueness theorem used in A.1
then puts $P$ in the stabilizer of the resulting primitive valuation.
Here, however, no point or cusp supplied a horizontal valuation in
advance. A valuation on a reduction component can have a nontrivial
geometric stabilizer. Thus this does not identify a cusp.
At the primes above $\ell$, splitting over the cyclotomic kernel
is also insufficient: a section over that infinite-extension group
need not be a point or cusp section over the local number field.

The next concrete attack is to use the embedding of the same
$P\times C$ in all finite cover groups, together with these local
valuation stabilizers, to force the compatible extension classes
of $B$ to vanish. Neither a compatible family of such vanishings
nor a configuration embedded in $\Pi$ has been constructed here.
\end{proof}

\begin{statement}{A.7}{Open}{The remaining assertion in R27}
The claim that every configuration $(P,C,B)$ in the actual curve
group virtually splits remains open in this return. The residual
and local-kernel statements below, and the audited CM restrictions
and prime-complement construction, do not prove that claim.
The local section theorem of the supplied October 2026 preprint
does imply the cuspidality of a virtually split normalizer; this
is a different direction from the required virtual splitting.
The alleged density gap in that preprint's Section 8 has the
source mismatch described in A.20.
\end{statement}
\begin{proof}[Scope of the audit]
The proofs of LM6 A.36, A.50--A.53, A.61--A.65 were read in the
attached file \texttt{LM6\_POINTS\_AS\_SECTIONS\_paper.tex}.
The arithmetic commutant arguments A.54--A.60 were independently
checked, including fresh finite-field computations; their results
belong to Part B of this return. We do not treat a statement-only
intake report as evidence for a proof. The independent arguments
and the exact source hypotheses used in this audit are recorded
in the following statements and in the work report.
\end{proof}

\begin{statement}{A.15}{Proved}{The residual group in LM6 A.36}
Let $P_n=\pi'\pi^n$ be the geometric group of the multiplication
cover $E\setminus E[n]\to E\setminus\{O\}$. Inside $\pi'$,
\[
 R_{\rm nil}(\pi')
  =\bigcap_{n\ge1,\ \ell}\ker(P_n\longrightarrow P_n^{(\ell)}).
                                                               \tag{A.15.1}
\]
Any noncuspidal closed procyclic subgroup with open arithmetic
normalizer is contained in this group, which lies in $\pi''$
and is strictly smaller than $R_{\rm nil}(\pi)$.
\end{statement}
\begin{proof}
If the subgroup has nonzero image in some $\pi^{(\ell)}$, the
criterion for the CM curve puts that image in inertia. Its
nontrivial $\ell$-component maps injectively to this procyclic
image, so its full normalizer acts cyclotomically on that
component. Nakamura's full profinite recognition theorem
\cite[Theorem 2.1.1]{Nakamura1994} puts that component in actual
full inertia. The full centralizer of any nonidentity inertia
element is that inertia; hence the whole procyclic subgroup is
cuspidal. This argument uses the theorem for the full class of
finite groups, not a lifting of conjugacy from a pro-prime quotient.

In the remaining case $A\subset R_{\rm nil}(\pi)\subset\pi'$,
so $A\subset P_n$ for all $n$. Suppose it survives in
$P_n^{(\ell)}$. Its proper homology maps to zero under the map
to $\pi^{(\ell)}$. On the proper Tate modules that map is
multiplication by $n$, which is injective even if $\ell\mid n$.
Thus its proper homology was zero. By A.11 its nonzero pro-$\ell$
image must be cuspidal, and the same full recognition and
centralizer argument gives a contradiction. This proves the
containment in the right side of (A.15.1).

Every map $P_n\to P_n^{(\ell)}$ kills $R_{\rm nil}(\pi')$ on
restriction to $\pi'$. Conversely, if an element of $\pi'$ is
not in $R_{\rm nil}(\pi')$, it survives in a finite quotient
$G$ of $\pi/R_{\rm nil}(\pi')$. This quotient exists because
$R_{\rm nil}(\pi')$ is characteristic in $\pi'$ and normal in
$\pi$. The image $G'$ is nilpotent. Take $n$ divisible by the
exponent of $G/G'$. The image of $P_n$ is exactly $G'$; a Sylow
projection of $G'$ detects the selected element. This proves
(A.15.1). Its containment in $\pi''$ follows because finite
abelian quotients are among the nilpotent quotients of $\pi'$.

For strictness use $G=\mathbb F_{25}\rtimes C_3$, with $C_3$
acting by a primitive cube root in $\mathbb F_{25}^{\times}$.
This is a two-generated quotient of $\pi$.
The executed program \texttt{a\_normalizer\_certificates.py}
checks its order $75$, $|G'|=25$, $G''=1$, and
$[G',G]=G'$. Thus its nilpotent residual is $G'\ne1$ while the
nilpotent residual of $G'$ is trivial. The image of a profinite
nilpotent residual in a finite quotient is the finite nilpotent
residual: use the descending closed lower-central series and
compactness to commute its intersection with the finite image.
Consequently the two residuals of $\pi$ in question are distinct.
\end{proof}

\begin{statement}{A.16}{Proved}{The CM and prime-support restrictions in LM6 A.52--A.53}
Let $A$ be a noncuspidal procyclic subgroup of the full geometric
group of a number-field hyperbolic curve, with normalizer image
$G_L$. Then $L$ contains a CM subfield. If $A_\ell\ne1$ and
$M/\Q$ is finite Galois containing $L$, some CM subfield of $M$
embeds into $\Q_\ell$. In particular, for imaginary quadratic
$L$ such a prime must split in $L$.
If $\ell$ is unramified in $M$ and its Frobenius conjugacy class
is that of a complex conjugation, then $A_\ell=1$. These excluded
primes have positive Dirichlet density. Thus a subgroup with
cofinite prime support, in particular one isomorphic to
$\widehat\Z$, must be cuspidal. The same is true with arbitrary
prime support if $L$ is totally real.
\end{statement}
\begin{proof}
We include the character descent that is essential in the full
group. Put $D=N_\Pi(A)\cap\pi$. For a nonzero component
$A_\ell\simeq\Z_\ell$, let $\theta$ denote its normalizer
character. Every pro-$\ell$ subgroup of $D$ containing $A_\ell$
is free pro-$\ell$ in the affine case. Its conjugation character
on $A_\ell$ is trivial: the first nonzero Magnus term of a
generator is unchanged by an inner automorphism, and cancellation
in its torsion-free graded piece forces exponent one. It follows
that the pro-$\ell$ Sylow subgroup of $\theta(D)$ is trivial.
Since $\Z_\ell^\times$ has prime-to-$\ell$ torsion of exponent
$\ell-1$ for odd $\ell$, and is pro-$2$ at two,
\[
 \lambda=\theta^d:G_L\longrightarrow\Z_\ell^\times,
 \qquad d=\ell-1\ (\ell>2),\quad d=1\ (\ell=2)
                                                               \tag{A.16.1}
\]
is a well-defined continuous character. For proper curves the
same argument is applied to each finitely generated pro-$\ell$
subgroup, using A.23 below.

Detect a generator of $A_\ell$ in proper homology of a finite
cover, using Nakamura's Lemma 2.1.4. Such a detector exists
because otherwise $A_\ell$, and hence $A$, would be cuspidal.
Intersect its arithmetic group with the normalizer and replace
the number field by the resulting finite extension $L'/L$.
The detected rational line has a character $\psi$ with
\[
                    \lambda|_{G_{L'}}=\psi^d.
                                                               \tag{A.16.2}
\]
It has Frobenius weight one, and its Hodge--Tate exponent relative
to the cyclotomic character is either zero or one at each
embedding. Hodge--Tate representations descend through finite
extensions, so $\lambda$ itself is Hodge--Tate; see
\cite[Chapter III, Appendix, Corollary 2]{A_Serre}.
The one-dimensional local algebraicity theorem and the global
unit relations in \cite[Chapter III, \S1.2; Chapter II, \S3.1]{A_Serre}
therefore apply. The character is unramified outside finitely
many primes, as is seen after the finite extension from the
Jacobian representation.

Write $G=\Gal(M/\Q)$, $H=\Gal(M/L)$. The algebraic type of
$\lambda|_{G_M}$ is an integer-valued function $n:G\to\Z$,
right $H$-invariant by norm pullback. The global unit relation
says that, with $a=|G|^{-1}\sum_g n(g)$ and $b=n-a$,
\[
                b(cg)=-b(g)
       \quad\text{for every complex conjugation }c.
                                                               \tag{A.16.3}
\]
The function $b$ cannot be zero. If it were, $a$ would be an
integer and the character would be finite times $\chi_\ell^a$:
the quotient has trivial algebraic type, is trivial on open
local units at $\ell$, and factors through a finite ray class
group. Kill the finite character and compare (A.16.2).
Frobenius absolute values give $2a=d$, whereas Hodge--Tate
exponents give $a\in\{0,d\}$, impossible for $d>0$.

Let $J\lhd G$ be generated by all products of two complex
conjugations. Equation (A.16.3) makes $b$ left $J$-invariant,
hence also right $J$-invariant by normality. It is therefore
right $JH$-invariant. Modulo $J$, all complex conjugations have
the same central image. None lies in $JH$, since it would then
both fix and negate the nonzero function $b$. Consequently
$M^{JH}\subset L$ is totally imaginary with a common
complex-conjugation involution and totally real fixed field;
it is a CM field.

For the local assertion take the left stabilizer
$P=\operatorname{Stab}_G(b)$. It contains $J$ and contains no
complex conjugation, so the same argument makes $M^P$ CM.
The algebraic type is defined over $\Q_\ell$; hence the local
decomposition subgroup $D_v\subset G$ fixes $b$ and lies in
$P$. In the selected completion, $M^P$ therefore embeds in
$M_v^{D_v}=\Q_\ell$. This includes ramified primes. If $L$ is
imaginary quadratic, take $M=L$; it embeds in $\Q_\ell$ exactly
at its split primes.

If Frobenius at an unramified $\ell$ is complex conjugation,
choose its prolongation so $D_v$ contains that conjugation.
Then (A.16.3) and $D_v$-invariance force $b=0$, contradiction.
Chebotarev gives density $|\operatorname{Cl}(c)|/|G|>0$.
Finally a totally real field contains no CM subfield. These
observations prove all the remaining assertions.
\end{proof}

\begin{statement}{A.17}{Proved}{The prime-complement configuration and the direction supplied by the local theorem}
For a noncuspidal normalizer in the affine setting there are,
after replacing it by an open subgroup and the field by a
finite extension, closed subgroups satisfying
\[
 \begin{gathered}
 N_0\cap\pi=P\times C,\qquad P\simeq\Z_\ell,
 \qquad 1\ne P\subset A_\ell,\qquad C\text{ pro-prime-to-}\ell,\\
 N_0=P\rtimes B,\qquad 1\to C\to B\to G_L\to1.
 \end{gathered}                                         \tag{A.17.1}
\]
The construction of (A.17.1) uses no section conjecture.
Using Theorem 1.1 of the supplied local section-conjecture
preprint as a professional-preprint input, a procyclic
normalizer is virtually split if and only if it is cuspidal.
Thus in (A.17.1) $C\ne1$ and $B\to G_L$ has no section
after any finite extension. This verifies precisely the two
different dependency levels of LM6 A.61--A.62.
\end{statement}
\begin{proof}
First shrink the normalizer so that its geometric kernel $D$
centralizes $A_\ell$, using the finite character image in
(A.16.1). A pro-$\ell$ Sylow subgroup $P^*$ of $D$ contains
the normal subgroup $A_\ell$ and is free pro-$\ell$. Its
nontrivial center makes it procyclic; thus $A_\ell$ is open
in $P^*$. The compact set $P^*\setminus A_\ell$ avoids one.
Choose an open normal subgroup $N_1$ of the normalizer avoiding
this set. Then $P_1=P^*\cap N_1$ is a nonzero central Sylow
subgroup of $D_1=D\cap N_1$, so it is the unique Sylow subgroup.
At every finite quotient a central Sylow subgroup has a unique
prime-to-$\ell$ complement: coprime cohomology gives existence,
and the absence of homomorphisms into the Sylow subgroup gives
uniqueness. These complements are compatible. Hence
$D_1=P_1\times C_1$, and uniqueness makes $C_1$ normal in $N_1$.

Choose a finite cover detecting $P_1$ in homology and intersect
its arithmetic group with $N_1$. Closed subgroups of a product
of a pro-$\ell$ and a pro-prime-to-$\ell$ group split into
their primary factors, by taking limits of integer powers of
each element. We retain the form $D=P\times C$ with $P\ne1$.
Its homology image must have nonzero proper projection:
otherwise it would be cyclotomic and full inertia recognition
would make $A$ cuspidal. Thus $P$ injects equivariantly into
$T_\ell J$ for the compactified cover's Jacobian.
Tate semisimplicity gives a rational equivariant retraction.
Clearing denominators gives $r:T_\ell J\to P$ with
$r|_P=\ell^k$ for some $k\ge0$.

Make a point of the compactified cover rational and use its
Abel--Jacobi map to the actual extension
$(\prod_qT_qJ)\rtimes G_L$. Restriction to $N$ and composition
with $r$ kills $C$ and embeds
\[
                   N/C\hookrightarrow P\rtimes G_L.
\]
Its image meets $P$ in $\ell^kP$ and surjects to $G_L$.
Modulo $\ell^kP$ it is the graph of a finite continuous
cocycle. Its zero set is an open subgroup. After restricting
there, subtracting the geometric coordinate of each lift gives
all $(0,g)$ in the image. Compactness makes the inverse onto
this closed image continuous, so $N/C$ has an actual section.
Let $B$ be its inverse image in $N$. Then $N=P\rtimes B$,
$B\cap\pi=C$, and (A.17.1) follows.

We give a shorter independent proof of the final dependency
than the infinite-index Hom argument in LM6 A.61. Suppose an
actual section $s:G_L\to N_\Pi(A)$ exists. Localize at one
finite place. The supplied local theorem
\cite[Theorem 1.1]{A_OpenAILocal}, with its affine neighborhood
consequence, makes this section a point or cusp section.
Here is that consequence: use a cofinal system of characteristic
open geometric subgroups $N$ and the actual open arithmetic
subgroups $N\rtimes s(G_{L_v})$. These are section neighborhoods
over the same local field. One can first choose a member whose
compactification has genus at least two: take a characteristic
cover in which every cusp inertia has sufficiently large order,
and apply Riemann--Hurwitz. Apply the proper theorem to every
further member of this cofinal system. The
resulting compactification points are compatible, by injectivity
in the proper theorem. Their finite fibres above the first
point have nonempty compatible inverse limit. The section
fixes this lifted horizontal point; its stabilizer is the
ordinary decomposition group or the cuspidal decomposition
group according to whether the point is in the affine curve.
This uses full etale covers and does not replace them by
pro-prime covers.

Conjugation on the procyclic group $A$ has abelian automorphism
group. Thus the actual section's restriction to the local
commutator centralizes $A$. The local assertion proved in A.1
now gives $A=1$ for an ordinary point section, or $A\subset I$
for a cusp section. The former is excluded. Conversely a
rational tangent at a rational cusp gives a section normalizing
every closed subgroup of its inertia, after finite extension.
This proves the equivalence. Applying it to (A.17.1) proves
the stated failure of virtual splitting and $C\ne1$.
Only this last paragraph and its affine-neighborhood input
depend on the supplied local section theorem; none uses the
global density inference discussed in A.20.
\end{proof}

\begin{statement}{A.18}{Proved}{Why the abstract kernel and local splitting data do not suffice}
For every prime $\ell$ and an odd prime $q\ne\ell$, there is
an extension of $G_{\Q(i)}$ by $\Z_q(1)$ with no section over
any finite extension. It can be chosen split at every place
above $\ell$. For a split $\ell$, adjoining the semidirect
factor $\Z_\ell(\psi)$ for a CM character gives an abstract
configuration with geometric kernel $\Z_\ell\times\Z_q(1)$,
which itself embeds as a closed subgroup of the free profinite
geometric group. No embedding of the whole arithmetic
configuration in $\Pi$ is asserted. This proves the obstruction
claimed in LM6 A.63 and isolates the additional embedding issue.
\end{statement}
\begin{proof}
Choose two finite places $v,w$ away from $\ell$ and $q$.
The Brauer invariant-sum sequence supplies compatible classes
$\beta_n\in H^2(G_K,\mu_{q^n})$ with local invariants
$q^{-n}$ and $-q^{-n}$ at $v,w$, and zero elsewhere.
The coefficient map $\mu_{q^{n+1}}\to\mu_{q^n}$ corresponds
to multiplication by $q$ on Brauer classes, so these classes
are compatible. Starting with a continuous normalized cocycle
representing $\beta_n$, choose one for $\beta_{n+1}$; their
difference after reduction is a coboundary of a finite-valued
continuous one-cochain. Lift that cochain value by value and
subtract its coboundary. Induction gives compatible actual
cocycles, whose inverse limit is a continuous
$z:G_K^2\to\Z_q(1)$.

On the product space define
\[
              (a,g)(b,h)=(a+g b+z(g,h),gh).
                                                               \tag{A.18.1}
\]
The cocycle identity gives a profinite group $B$. At a place
$v'$ above $v$ of a finite extension $L/K$, restriction of
$\beta_n$ has invariant $[L_{v'}:K_v]/q^n$, nonzero for
sufficiently large $n$. A section of $B\to G_K$ over $G_L$
would make all these finite reductions coboundaries, a
contradiction. This verifies continuous nonsplitting directly;
it does not infer a continuous obstruction from an unspecified
inverse limit of vanishing finite obstructions.

At a place above $\ell$ all the finite local classes are zero.
The local finite groups $H^1(G_{K_v},\mu_{q^n})$ satisfy
Mittag--Leffler, so the continuous comparison has no
$\varprojlim{}^1$ kernel. Thus the continuous local class is
zero and the extension splits there.
If $\ell$ splits in $K$, choose one CM Tate character $\psi$
of $E$ and put $N=\Z_\ell(\psi)\rtimes B$, with $B$ acting
through its quotient $G_K$. Its geometric kernel is the direct
product of the two different prime factors. Abstractly this is
a closed procyclic subgroup of $\widehat\Z$, hence embeds in
a free profinite generator subgroup and is projective. The
action, centralization and local splitting properties therefore
do not alone rule out the nonsplitting. The missing condition
is an embedding of this entire extension, with its arithmetic
projection and action, in the particular curve group $\Pi$.
\end{proof}

\begin{statement}{A.23}{Proved}{The proper-curve version of the audit}
Every finitely generated pro-$\ell$ subgroup of the full
geometric group of a proper hyperbolic curve is free pro-$\ell$.
The CM and prime-support restrictions of A.16 and the
prime-complement configuration of A.17 consequently hold
also for proper hyperbolic curves. There is no boundary
alternative. Using the supplied local section theorem, their
normalizers cannot virtually split. These are the assertions
of LM6 A.64--A.65 with the dependencies separated.
\end{statement}
\begin{proof}
Let $Q$ be such a finitely generated pro-$\ell$ subgroup of
a proper surface group $\Pi_g$, $g\ge2$.
Continuous finite-coefficient cohomology of $Q$ is the direct
limit of that of open overgroups $V$ in $\Pi_g$.
The subgroup $Q$ cannot be open: every open surface group has
nontrivial abelian quotients at primes other than $\ell$.
If the supernatural index $[\Pi_g:Q]$ had finite $\ell$-adic
valuation, choose an open overgroup $V$ for which the
$\ell$-part of $[\Pi_g:V]$ is maximal, equivalently the
remaining $\ell$-part of $[V:Q]$ is minimal. All further
restriction maps toward $Q$ then have prime-to-$\ell$ degree
on a cofinal subsystem and are injective on $H^1(-,\F_\ell)$
by corestriction. The genera of these covers are unbounded,
so $H^1(Q,\F_\ell)$ would have unbounded dimension, contrary
to finite generation. Thus $[\Pi_g:Q]$, and hence every
$[V:Q]$, has infinite $\ell$-adic valuation.
On $H^2(V,\F_\ell)$ restriction is
multiplication by the covering degree, hence is eventually
zero. Thus $H^2(Q,\F_\ell)=0$, which is the free pro-$\ell$
criterion.

For the Sylow step of A.17 no finite generation of the whole
Sylow group is needed. Each subgroup generated by $A_\ell$
and one additional Sylow element is finitely generated,
free pro-$\ell$ by the preceding argument, and has central
nonzero $A_\ell$ after the character shrinking. It is cyclic.
Every finite image of the whole Sylow group is therefore
cyclic, so the Sylow group is procyclic. The rest of the
prime-complement and semisimple-retraction construction is
unchanged, with proper homology as the detector from the
outset. Finally a local section of a proper curve is an
ordinary point section by the supplied theorem. Its local
derived centralizer is trivial by A.1, contradicting any
nonzero normalized procyclic subgroup. This proves the last
assertion with exactly the stated preprint dependency.
\end{proof}

\begin{statement}{A.19}{Proved}{Exactly what the two local kernels differ by}
Let $v\mid p$ be a finite place of $K=\Q(i)$ and put $k=K_v$.
Under a compatible decomposition-group embedding,
\[
 D_v\cap\Delta=\overline{[G_k,G_k]},\qquad
 \frac{G_{k(\mu_\infty)}}{\overline{[G_k,G_k]}}
 \simeq\ker\bigl(N_{k/\Q_p}:\mathcal O_k^{\times}
                                  \to\Z_p^{\times}\bigr).
                                                               \tag{A.19.1}
\]
The second group is trivial at split primes, is noncanonically
$\mu_{p+1}\times\Z_p$ at odd inert primes, and is
$\mu_4\times\Z_2$ at the ramified prime two.
For a finite extension $L/K$ and a place $w\mid v$, the kernel
using the original $\Delta$ is
$G_{L_w}\cap\overline{[G_k,G_k]}$; the quotient lost by passing
from $G_{L_w(\mu_\infty)}$ to this kernel is an open subgroup
of the norm kernel in (A.19.1).
\end{statement}
\begin{proof}
The easy containment of the local commutator in $\Delta$ follows
from abelianization. For equality, every finite abelian extension
of $k$ occurs as the completion of a finite abelian extension of
$K$: the Grunwald--Wang theorem has no exceptional condition here
because $i\in K$. See \cite[Theorems 1--2]{A_LorenzRoquette}.
Thus the map $G_k^{\rm ab}\to G_K^{\rm ab}$ is injective,
which gives the first equality. The same local decomposition
identification is stated in \cite[Chapter III, \S2.1]{A_Serre}.

Local reciprocity identifies $G_k^{\rm ab}$ with the profinite
completion of $k^{\times}$. Its action on all roots of unity
is the restriction map to $G_{\Q_p}^{\rm ab}$, identified by
local Kronecker--Weber with all the local cyclotomic characters.
Reciprocity is compatible with norms. Hence the quotient in
(A.19.1) is the kernel of the completed norm
$\widehat{k^{\times}}\to\widehat{\Q_p^{\times}}$.
On valuation factors the norm is multiplication by the residue
degree, injective on $\widehat\Z$. Its kernel is therefore
exactly the displayed unit norm kernel, with no valuation factor.

For $k=\Q_p$ it is zero. For the unramified quadratic extension
at odd $p$, the norm on residue roots of unity has kernel
$\mu_{p+1}$, and on principal units the logarithm identifies
the norm with trace on a rank-two $\Z_p$-lattice, giving a
rank-one kernel. For $k=\Q_2(i)$ the kernel has torsion
exactly $\mu_4$, and its principal-unit part has rank one,
again by the logarithm on a sufficiently small open subgroup.
The classification of finitely generated abelian pro-$2$ groups
gives the stated noncanonical product.

Finally the chosen local embedding for $L_w$ is an open subgroup
of $G_k$. Intersecting the first equality in (A.19.1) gives the
claim about the original $\Delta$. The image of an open subgroup
under a surjective homomorphism of profinite groups is open,
which proves the last assertion. If one instead replaces
$\Delta$ by the new global commutator of $G_L$, its local
intersection is $G_{L_w}'$ and the principal norm-kernel rank
is $[L_w:\Q_p]-1$; these are different comparisons.

The loss is consequently explicit. In A.1 it loses no uniqueness
hypothesis: the smaller local derived subgroup already has a
nonzero tame $\Z_r$ quotient for every $r\ne p$. For an arbitrary
normalizer, replacing the larger local kernel by the smaller one
still does not supply a point or cusp section to which the
centralizer theorem could be applied.
\end{proof}

\begin{statement}{A.20}{Refuted}{The density intersection inference, and its correct attribution}
The inference that a set of strictly positive Dirichlet density
must meet a set of potentially density one is false. The
counterexample below refutes that inference only. It does not
refute a section theorem or a curve finiteness theorem.
The attached local section-conjecture preprint's Section 8
does not contain the alleged inference or a Corollary 7.10.
LM6 D.6 attributes that step to the other 2026 source, Porowski's
unnumbered corollary following Corollary 7.10.
\end{statement}
\begin{proof}
Let $L/K$ be a nontrivial quadratic extension. The unramified
inert primes of $K$ have density $1/2$. The completely split
primes also have density $1/2$, and are disjoint from them.
But the pullback of the completely split set to $L$ has density
one: primes over inert primes have relative residue degree two
and contribute a bounded series at the Dirichlet-density pole.
Thus the split set has potentially density one and does not meet
the positive-density inert set.

Replacing strictly positive by strongly positive density repairs
this particular inference: by definition its pullback has
positive upper density after every finite extension, so it
meets any set that becomes density one in some such extension.
This is exactly the hypothesis retained by the valid uniform
reduction argument in LM6 D.6. No counterexample to the stronger
arithmetic conclusion is inferred from the failed weaker proof.

We checked the whole supplied file
\texttt{sections/08-global-sections.tex}. It proceeds by
Harari--Stix finite descent, Stix's finite-support result and
Stoll's finite-descent criteria; it explicitly retains the
finite-cover equality as a hypothesis. It has no use of the
displayed false intersection step. The local Theorem 1.1 is
proved before that section. Consequently this source mismatch
does not invalidate that local input. The arguments A.1, A.3,
A.6, A.11--A.15, and A.19 do not use either the local section
theorem or the disputed global density inference.
\end{proof}
